\section{Part II: Self-Consistent Calibration}
\label{sec:part2}

\subsection{The calibration fixed point}

Activation-aware post-training quantization chooses a code for layer $\ell$ by
minimizing a metric weighted by that layer's input statistics
$H^{xx}_\ell=\mathbb E[x_\ell x_\ell^\top]$, estimated on calibration data.
GPTQ and AWQ measure those statistics on the full-precision model
\cite{gptq,awq}. Once many layers are replaced, the deployed model presents
layer $\ell$ with a different distribution from the one its codebook was fitted
to, so the conventional procedure computes
$Q_{\text{one-shot}}=\mathcal Q(H[W_{\rm BF16}])$ while the quantity describing
the deployed model is the fixed point
\begin{equation}
 Q^\star=\mathcal Q\bigl(H[Q^\star]\bigr).
 \label{eq:fp}
\end{equation}

Equation~\eqref{eq:fp} has the form of the Kohn--Sham self-consistency condition
\cite{kohnsham}, which suggested treating it the same way: statistics in place
of density, calibration loss in place of total energy, linear or Anderson
mixing \cite{anderson,pulay} in place of density mixing, and codebook-assignment
churn in place of density convergence. We implemented that loop, with a guard
that rejects any iterate raising the calibration loss and halves the mixing
parameter.

This section reports that the analogy is wrong in the specific way that matters,
and that the repair we designed for it also fails.

\subsection{Why one pass suffices: acyclicity}
\label{sec:acyclic}

The statistics $H^{xx}_\ell$ depend only on layers $1,\dots,\ell-1$. Quantizing
in topological order while propagating the quantized activations forward
therefore gives every layer exactly the statistics it will see in the deployed
model. The dependency graph is acyclic, the fixed point of \eqref{eq:fp} is
reached in \textbf{one sweep}, and there is nothing for an iteration to do.

This is not merely a theoretical point. Two published methods already exploit
it: GPTAQ's asymmetric calibration \cite{gptaq} and CoreQ's mismatch correction
\cite{coreq} both target accumulated error from previously quantized layers, and
both are closed-form and single-pass. We should have derived the acyclicity
argument before building an iteration, and we state it here because it is the
reason the iteration cannot be the contribution.

\paragraph{Measured.} Table~\ref{tab:phase0a} compares the one-pass sequential
arm against our iterative loop on the 2-D codec, the codec most sensitive to the
statistics. The prediction was redundancy; the measurement is stronger. The
one-pass arm beats the iteration by $0.004228$ NLL and beats one-shot
calibration by $0.002778$.

\begin{table}[ht]
\centering\small
\caption{Forward-only self-consistency. Qwen3.5-0.8B, 2-D codec, all 18 DeltaNet
input projections, 32{,}768 WikiText-2 validation targets.}
\label{tab:phase0a}
\begin{tabular}{lrrl}\toprule
Calibration & NLL & $\Delta$NLL & 95\% CI\\\midrule
BF16 reference & 3.233953 & --- & ---\\
one-shot (BF16 activations) & 3.245136 & $+0.011184$ & $[+0.009132,+0.013299]$\\
\textbf{sequential, one pass} & \textbf{3.242358} & $\mathbf{+0.008406}$ & $[+0.006635,+0.010215]$\\
iterative SCF loop & 3.246586 & $+0.012633$ & $[+0.010697,+0.014610]$\\\bottomrule
\end{tabular}
\end{table}

The mechanism is clear in hindsight. Sequential calibration uses statistics that
are \emph{correct}, not approximate. Our loop mixes a \emph{global} statistics
vector toward $H[Q]$ with $\alpha<1$; the mixed statistic corresponds to no
actual model state -- it interpolates between two different networks -- and all
eighteen layers are refitted against it simultaneously. Damping is what makes
SCF stable in density functional theory; here it is what makes it wrong, because
the exact answer is available in one sweep and damping moves away from it.

\paragraph{But solving the forward fixed point correctly does not help much
either.} Table~\ref{tab:phase0a} is a 32{,}768-target screening slice. Repeating
the comparison on the complete validation stream removes the advantage:
sequential calibration changes the Lloyd scalar code by $+0.000022$, leaves
decoupled polar \emph{bit-identical} (that codec ignores activation statistics
entirely), improves repaired polar by $-0.000596$, and \emph{worsens} the 2-D
codec by $+0.000607$ (Table~\ref{tab:main}). The effects are small and change
sign across codecs. The iterative loop behaves the same way across calibration
sizes: on the full stream it worsened the 2-D codec by $+0.002720$ with sixteen
shared calibration windows and improved it by $-0.000585$ with sixty-four
windows and a disjoint guard set.

The explanation is in the residual itself. We replace one tensor in each of the
eighteen DeltaNet layers of this 24-layer model -- $15.1\%$ of its parameters --
and the forward statistics move by only $0.87\%$, so there is very little for any
forward-only correction to recover. Sequential propagation is the
\emph{correct} way to solve that fixed point -- the acyclicity argument is
unaffected -- but at this scope the fixed point is nearly where one-shot
calibration already sits. We would expect the correction to matter more under
whole-model quantization, which we did not run.

We also note the methodological point: the screening slice favoured sequential
calibration by $0.002778$ NLL and the full stream did not reproduce it. Reduced
evaluation slices in this project have repeatedly produced effects that vanish
at full scale, and we treat the screening result as superseded.

\subsection{The bidirectional reformulation}

The fixed point becomes non-trivial exactly when the per-layer objective depends
on layers \emph{downstream}. Expanding the calibration loss to second order in
each layer's output perturbation $\Delta y_\ell=\Delta W_\ell x_\ell$, with the
Gauss--Newton/Fisher approximation to the output Hessian,
\begin{equation}
 \text{cost}_\ell=\sum_j S_{\ell,j}\,\Delta w_{\ell,j}^\top H^{xx}_\ell \Delta w_{\ell,j},
 \qquad
 S_{\ell,j}=\mathbb E\!\left[\left(\frac{\partial\mathcal L}{\partial y_{\ell,j}}\right)^{\!2}\right].
\end{equation}
Here $H^{xx}_\ell$ depends on upstream layers and $S_{\ell,j}$ on downstream
ones, because the gradient at layer $\ell$ is backpropagated through every later
layer and those are quantized too. The dependency graph is now \textbf{cyclic}:
no topological order exists, so no single sweep in either direction satisfies
both conditions. This is the Kohn--Sham situation proper -- the potential
depends on the density everywhere, not only upstream -- and it is the part that
\cite{gptaq,coreq} do not address, because a layer-local output-matching target
assumes all output errors matter equally, which is what $S_\ell$ denies.

\paragraph{The coupling is large.} We measured the relative residual of each
statistic when the eighteen projections are quantized. The forward residual is
$0.87\%$ pooled and grows monotonically with depth ($0.336\%$ at layer 1 to
$1.582\%$ at layer 22), consistent with accumulated upstream error. The backward
residual is $7.33\%$ pooled, $8.41\times$ larger, and is not monotone in depth.
Layer 0 is the clean illustration: its forward residual is \emph{identically
zero}, because its input is upstream of every quantized layer, while its
backward residual is $6.95\%$ because everything downstream of it is quantized.

So the coupling the reformulation identifies is real and an order of magnitude
larger than the one the literature corrects. That made the next test worth
running, and it is where the direction dies.

\subsection{Fisher weighting does not rank layers}
\label{sec:fisherfail}

If $S$-weighted error is the right local surrogate for end-to-end damage, it
should predict which layers are costly to quantize. We quantized one layer at a
time, all others BF16, and correlated four predictors against the measured
validation damage over the eighteen layers.

\begin{table}[ht]
\centering\small
\caption{Predicting single-layer quantization damage. 65{,}536 validation
targets; 10 of 18 layers have a $\Delta$NLL interval excluding zero.}
\label{tab:fisher}
\begin{tabular}{lrrr}\toprule
Predictor & Spearman & Pearson & Pearson without L0\\\midrule
Weight MSE & $-0.218$ & $-0.202$ & $+0.072$\\
Diagonal activation-weighted & $+0.020$ & $-0.061$ & $+0.146$\\
Exact output error (forward only) & $+0.005$ & $-0.057$ & $+0.145$\\
Fisher (forward $\times$ backward) & $+0.015$ & $+0.898$ & $+0.181$\\\bottomrule
\end{tabular}
\end{table}

The Fisher Pearson coefficient of $0.898$ is an artifact of a single leverage
point. Excluding layer 0 it collapses to $0.181$, and its Spearman coefficient
is $-0.169$. The disagreement between a high Pearson and a null Spearman is the
diagnostic: no predictor, Fisher included, has rank skill across the other
seventeen layers.

Two limits on this test's power should be stated. Eight of eighteen layers have
damage intervals covering zero, so the comparison among low-damage layers is
weak. And three layers have \emph{significantly negative} damage -- quantizing
them improves validation NLL -- which no unsigned error surrogate can predict at
all. That caps the achievable correlation independently of the metric.

\subsection{The one real finding, and why it is not exploitable}

Layer 0 is simultaneously the most damaging layer to quantize, the layer Fisher
ranks highest, and the layer weight MSE ranks \emph{lowest} -- that is, the
criterion used throughout the earlier phases of this work identifies the single
most damaging tensor in the model as the safest one. Layer 0 alone accounts for
$58.2\%$ of the summed single-layer damage. This is a crisp mechanism for the
observation, recurring throughout this project, that weight MSE fails to predict
task loss.

It does not convert into a method. Table~\ref{tab:protect} asks which single
tensor is worth leaving in BF16. Fisher's choice recovers $57.7\%$ of the damage
and weight MSE's choice recovers $3.1\%$, which looks decisive until the right
control is added. Leaving one tensor unquantized costs $1.167\times$ the
all-quantized payload, so the honest comparison is a larger codebook applied
uniformly at the same budget.

\begin{table}[ht]
\centering\small
\caption{Protecting one tensor versus spending the same bytes uniformly.
Complete validation stream, 261{,}248 targets.}
\label{tab:protect}
\begin{tabular}{lrlrr}\toprule
Arm & $\Delta$NLL & 95\% CI & Recovered & Payload\\\midrule
All 18 quantized & $+0.008393$ & $[+0.007720,+0.009069]$ & --- & $1.000\times$\\
Protect L0 (Fisher's pick) & $+0.003716$ & $[+0.003197,+0.004245]$ & $55.7\%$ & $1.167\times$\\
Protect L20 (weight MSE's pick) & $+0.008061$ & $[+0.007401,+0.008729]$ & $4.0\%$ & $1.167\times$\\
Protect L10 (arbitrary) & $+0.008382$ & $[+0.007733,+0.009025]$ & $0.1\%$ & $1.167\times$\\
\textbf{All 18, 4.5-bit codebook} & $+0.004872$ & $[+0.004348,+0.005402]$ & $41.9\%$ & $1.125\times$\\
\textbf{All 18, 5-bit codebook} & $\mathbf{+0.002950}$ & $[+0.002555,+0.003346]$ & $\mathbf{64.9\%}$ & $1.250\times$\\\bottomrule
\end{tabular}
\end{table}

Interpolating the uniform-bits frontier to the same $1.167\times$ payload gives
$49.6\%$ recovery. Protecting the Fisher-identified layer is therefore worth
$+6.1$ percentage points, against a confidence half-width of $6.2$ points on the
protected arm alone, before accounting for the interpolation's own uncertainty.
\textbf{The advantage is at best marginal and we do not regard it as
established: outlier protection does not clearly beat spending the same bytes
uniformly.}

\subsection{Summary of Part II}

The forward-only self-consistency loop is redundant by acyclicity, already
addressed in closed form by published one-pass methods, and measurably worse
than the one-pass baseline. The bidirectional reformulation identifies a
genuinely cyclic coupling that is $8.41\times$ larger than the one the
literature corrects, but the metric it implies has no rank skill across layers,
and its single correct call yields no advantage over uniform bit allocation at
matched bytes. We do not recommend pursuing the direction further.
